\subsection{Continuity}

Next goal: study functions $f:U\to\R$, where $U\subseteq\R$,
usually $U=\R$ or $(a,b)$ or $[a,b]$.

\begin{definition}{Continuity at a Point}{continuity-at-point}
  Let $U\subseteq\R$, $f:U\to\R$, and $x_0\in U$.
  \[
    \begin{aligned}
      f\text{ is continuous at }x_0
      \iff{} & \forall\varepsilon>0,\ \exists\delta>0\text{ s.t.} \\
             & \forall x\in U,\quad \abs{x-x_0}<\delta
      \implies\abs{f(x)-f(x_0)}<\varepsilon.
    \end{aligned}
  \]
\end{definition}

\begin{definition}{Continuity on a Set}{continuity-on-set}
  $f$ is continuous on $U$ $\iff$ $\forall x_0\in U$, $f$ is continuous at $x_0$.
\end{definition}

\begin{thinking}
  \begin{minipage}[c]{0.47\linewidth}
    \centering
    \begin{tikzpicture}[font=\small,declare function={f(\x)=0.4+0.23*(\x+0.3)^2;}]
      \draw[->] (-0.4,0) -- (3.6,0) node[right] {$x$};
      \draw[->] (0,-0.1) -- (0,3.6) node[above] {$y=f(x)$};
      \draw[sketch accent,domain=-0.35:3.3,samples=40] plot (\x,{f(\x)});
      \foreach \x/\lab/\pos in {1.7/{x_0-\delta}/below left,2.2/{x_0}/below,2.7/{x_0+\delta}/below right}{
        \draw[densely dashed] (\x,0) node[\pos] {$\lab$} -- (\x,{f(\x)});
      }
      \foreach \dy/\lab in {-0.7/{f(x_0)-\varepsilon},0/{f(x_0)},0.7/{f(x_0)+\varepsilon}}{
        \draw[densely dashed] (0,{f(2.2)+\dy}) node[left] {$\lab$} -- (2.85,{f(2.2)+\dy});
        \fill (0,{f(2.2)+\dy}) circle (1.2pt);
      }
      \node[sketch point] at (2.2,{f(2.2)}) {};
    \end{tikzpicture}
  \end{minipage}
  \hfill
  \begin{minipage}[c]{0.51\linewidth}
    \begin{enumerate}[leftmargin=*]
      \item For a small $\varepsilon$-window,
      \item choose a $\delta$-window
      \item s.t. $f\bigl((x_0-\delta,x_0+\delta)\cap U\bigr)$\\
            $\subseteq(f(x_0)-\varepsilon,f(x_0)+\varepsilon)$.
    \end{enumerate}
    Usually find $\delta=\delta(\varepsilon,x_0)$.
  \end{minipage}
\end{thinking}

\begin{example}
  $f(x)=3x-1$. Prove $f$ continuous at $x_0=1$.

  \textit{Proof sketch.}
  For $\varepsilon>0$, want $\abs{x-1}$ small enough so that $\abs{f(x)-f(1)}<\varepsilon$.
  \[
    \abs{f(x)-f(x_0)}=\abs{f(x)-f(1)}=3\abs{x-1}.
  \]
  Want it $<\varepsilon$. Want $\abs{x-1}<\frac13\varepsilon$.

  \begin{proof}[Formal proof]
    $\forall\varepsilon>0$, let $\delta=\varepsilon/3$.\\
    Suppose $x\in\R$, $\abs{x-1}<\delta=\varepsilon/3$.
    \[
      \abs{f(x)-f(1)}=3\abs{x-1}<3\frac{\varepsilon}{3}=\varepsilon. \qedhere
    \]
  \end{proof}
\end{example}

\begin{example}[Step function]
  \[
    \begin{tikzpicture}[baseline=(current bounding box.center),font=\small]
      \draw[->] (-0.3,0) -- (3.6,0) node[right] {$x$};
      \draw[->] (0,-0.2) -- (0,2);
      \draw[densely dashed] (0,1.3) node[left] {$1$} -- (1.8,1.3);
      \draw[sketch accent] (-0.3,0) -- (1.8,0);
      \draw[sketch accent] (1.8,1.3) -- (3.35,1.3);
      \draw[sketch accent,fill=white] (1.8,0) circle (2pt) node[below=2pt,black] {$x_0$};
      \node[sketch point] at (1.8,1.3) {};
      \draw[<->] (0.32,0.65) -- (0.32,1.3) node[midway,right] {$\varepsilon$};
      \draw[<->] (0.32,1.3) -- (0.32,1.95) node[midway,right] {$\varepsilon$};
    \end{tikzpicture}
    \qquad
    u_{x_0}(x)=
    \begin{cases}
      1 & x\ge x_0, \\
      0 & x<x_0.
    \end{cases}
  \]
  Negation of continuity (Definition~\ref{def:continuity-at-point}):
  \[
    \begin{gathered}
      \exists\varepsilon>0\text{ s.t. }\forall\delta>0,\ \exists x\in U\\
      \text{s.t. }\abs{x-x_0}<\delta\text{ and }\abs{f(x)-f(x_0)}\ge\varepsilon.
    \end{gathered}
  \]
  At $x_0$, pick window $\varepsilon=\tfrac12$.
  $\forall\delta>0$, take $x=x_0-\delta/2$. Then
  $\abs{u_{x_0}(x)-u_{x_0}(x_0)}=1>\varepsilon$.
\end{example}

\begin{example}
  Let $U=\R$, $f:\R\to\R$, $f(x)=2x^2+1$. Prove $f$ continuous everywhere.

  \textit{Sketch.} Want $\abs{x-x_0}$ small enough to ensure $\abs{f(x)-f(x_0)}<\varepsilon$.
  \[
    \begin{aligned}
      \abs{f(x)-f(x_0)}
       & =\abs{2x^2+1-2x_0^2-1}                             \\
       & =2\abs{x^2-x_0^2}                                  \\
       & =2\underbrace{\abs{x-x_0}}_{\text{some control}}\,
      \underbrace{\abs{x+x_0}}_{\text{how?}}.
    \end{aligned}
  \]
  If we assume $\abs{x-x_0}<1$,
  \[
    \begin{aligned}
      \abs{x}             & =\abs{x-x_0+x_0}\le\abs{x-x_0}+\abs{x_0}<1+\abs{x_0},                  \\
      \implies\abs{x+x_0} & \le\abs{x}+\abs{x_0}<\underbrace{1+2\abs{x_0}}_{\text{some constant}}.
    \end{aligned}
  \]
  Some control of $\abs{x-x_0}$ can yield control of $\abs{x+x_0}$.

  \begin{proof}
    Let $x_0\in\R$, $\varepsilon>0$. Let
    $\displaystyle\delta=\min\left\{1,\frac{\varepsilon}{2(1+2\abs{x_0})}\right\}$.\\
    Let $x\in\R$ s.t. $\abs{x-x_0}<\delta$.
    \proofgoal{Remaining goal:}{$2\abs{x-x_0}\abs{x+x_0}<\varepsilon$.}
    The first bound $\delta\le1$ gives $\abs{x+x_0}<1+2\abs{x_0}$.
    Use the second bound on $\delta$ for $\abs{x-x_0}$:
    \[
      \abs{f(x)-f(x_0)}=2\abs{x-x_0}\abs{x+x_0}
      <2\frac{\varepsilon}{2(1+2\abs{x_0})}(1+2\abs{x_0})=\varepsilon. \qedhere
    \]
  \end{proof}
\end{example}

\subsection{Sequential Criterion for Continuity}

\begin{theorem}{Sequential Criterion for Continuity}{sequential-criterion-continuity}
  $f:U\to\R$, $x_0\in U$.
  \[
    \begin{aligned}
       & f\text{ continuous at }x_0                            \\
       & \iff\bigl(\forall\text{ seq. }(x_n)\text{ in }U,\quad
      x_n\to x_0\implies f(x_n)\to f(x_0)\bigr).
    \end{aligned}
  \]
\end{theorem}

\begin{thinking}
  $f$ continuous $\iff f$ “commutes” with limits of sequences.
  \begin{center}
    \begin{tikzpicture}[font=\small,x=1.35cm,declare function={f(\x)=0.38+0.21*(\x+0.5)^2;}]
      \draw[->] (-0.5,0) -- (4.4,0) node[right] {$x$};
      \draw[->] (0,-0.15) -- (0,3.8);
      \draw[sketch accent,domain=-0.5:3.5,samples=40] plot (\x,{f(\x)});
      \foreach \x/\i in {1.3/0,2/3,2.5/2,3/1}{
        \draw[densely dashed] (0,{f(\x)}) node[left] {$f(x_{\i})$}
          -- (\x,{f(\x)}) -- (\x,0) node[below] {$x_{\i}$};
      }
      \node[below] at (1.65,0) {$\cdots$};
    \end{tikzpicture}
  \end{center}
\end{thinking}

\begin{proof}[$(\Rightarrow)$]
  Suppose $f$ continuous at $x_0$.\\
  Let $(x_n)$ be a seq. in $U$ s.t. $x_n\to x_0$.
  \proofgoal{WTS:}{$f(x_n)\to f(x_0)$.}
  Fix $\varepsilon>0$.
  \proofgoal{By definition:}{$\exists N\in\N,\ \forall n>N,\ \abs{f(x_n)-f(x_0)}<\varepsilon$.}
  \begin{enumerate}[proofparts]
    \item For this $\varepsilon$, continuity at $x_0$ gives $\delta>0$ s.t. for every $x\in U$,
          \[
            \abs{x-x_0}<\delta\implies\abs{f(x)-f(x_0)}<\varepsilon.
          \]
    \item For this $\delta$, convergence of $(x_n)$ gives $N\in\N$ s.t.
          $n>N\implies\abs{x_n-x_0}<\delta$.
  \end{enumerate}
  \begin{center}
    \begin{tikzpicture}[font=\small,declare function={f(\x)=0.4+0.22*(\x+0.4)^2;}]
      \draw[->] (-0.6,0) -- (4.1,0) node[right] {$x$};
      \draw[->] (0,-0.1) -- (0,3.8);
      \draw[sketch accent,domain=-0.6:3.45,samples=40] plot (\x,{f(\x)});
      \foreach \x/\lab/\pos in {1.7/{x_0-\delta}/below left,2.2/{x_0}/below,2.7/{x_0+\delta}/below right}{
        \draw[densely dashed] (\x,0) node[\pos] {$\lab$} -- (\x,{f(\x)});
      }
      \foreach \dy in {-0.7,0,0.7}{
        \draw[densely dashed] (0,{f(2.2)+\dy}) -- (2.85,{f(2.2)+\dy});
      }
      \node[left] at (0,{f(2.2)}) {$f(x_0)$};
      \draw[->] (-0.25,{f(2.2)+0.23}) -- (-0.25,{f(2.2)+0.7}) node[above] {$+\varepsilon$};
      \draw[->] (-0.25,{f(2.2)-0.23}) -- (-0.25,{f(2.2)-0.7}) node[below] {$-\varepsilon$};
      \node[sketch point] at (2.2,{f(2.2)}) {};
    \end{tikzpicture}
  \end{center}
  Use this $N$ and fix $n>N$. Since $x_n\in U$, the two implications give
  \[
    n>N\implies\abs{x_n-x_0}<\delta
    \implies\abs{f(x_n)-f(x_0)}<\varepsilon. \qedhere
  \]
\end{proof}

\begin{proof}[$(\Leftarrow)$ by contrapositive]
  Suppose $f$ not cont. at $x_0$.
  \proofgoal{WTS:}{$\neg\bigl(\forall (x_n)\text{ in }U,\quad x_n\to x_0\implies f(x_n)\to f(x_0)\bigr)$.}
  \proofgoal{Equivalently:}{$\exists (x_n)\text{ in }U,\quad x_n\to x_0\text{ and }f(x_n)\not\to f(x_0)$.}
  $f$ not cont. at $x_0\iff$
  \[
    \exists\varepsilon>0,\ \forall\delta>0,\ \exists x\in U,\quad
    \abs{x-x_0}<\delta\text{ and }\abs{f(x)-f(x_0)}\ge\varepsilon.
  \]
  Fix such an $\varepsilon$. For each $n\in\N$, take $\delta=1/n$ and choose $x_n\in U$ s.t.
  \[
    \abs{x_n-x_0}<\frac1n\quad\text{and}\quad\abs{f(x_n)-f(x_0)}\ge\varepsilon.
  \]
  \begin{center}
    \begin{tikzpicture}[font=\small,x=1.3cm]
      \draw[->] (-0.5,0) -- (5,0) node[right] {$x$};
      \draw[->] (0,-0.15) -- (0,3.2);
      \foreach \x in {1.2,1.9,2.4,2.7,2.9}{
        \draw[densely dashed] (0,{0.45+0.3*\x}) -- (\x,{0.45+0.3*\x}) -- (\x,0);
      }
      \foreach \x/\lab in {1.2/{x_1},1.9/{x_2},2.35/{\cdots},2.9/{x_n}}{
        \node[below] at (\x,0) {$\lab$};
      }
      \node[left] at (0,0.81) {$f(x_1)$};
      \node[left] at (0,1.32) {$f(x_n)$};
      \draw[densely dashed] (0,2.4) node[left] {$f(x_0)$} -- (3.1,2.4);
      \draw[densely dashed] (3.1,1.38) -- (3.1,0) node[below right] {$x_0$};
      \draw[sketch accent] (-0.45,0.315) -- (3.1,1.38);
      \draw[sketch accent,fill=white] (3.1,1.38) circle (2pt);
      \draw[sketch accent] (3.1,2.4) -- (4.6,3.1);
      \node[sketch point] at (3.1,2.4) {};
    \end{tikzpicture}
  \end{center}
  Use this sequence $(x_n)$. It remains to verify both:
  \begin{enumerate}[proofsubgoals]
    \subgoal{$x_n\to x_0$.}
    Since $0\le\abs{x_n-x_0}<1/n$ and $1/n\to0$, the squeeze theorem applies.
    \subgoal{$f(x_n)\not\to f(x_0)$.}
    For the same $\varepsilon>0$ and any $N$, take $n=N+1>N$.
    Then $\abs{f(x_n)-f(x_0)}\ge\varepsilon$.\qedhere
  \end{enumerate}
\end{proof}

\newpage
\subsubsection{Some Applications}

\begin{theorem}{Arithmetic of Continuous Functions}{arithmetic-continuous-functions}
  Let $f,g:U\to\R$, $f,g$ cont. at $x_0\in U$.
  The following are cont. at $x_0$ (pointwise operations):
  \begin{enumerate}
    \item $f+g$.
    \item $f\cdot g$.
    \item $f/g$ on its domain if $g$ is nonzero at $x_0$.
  \end{enumerate}
\end{theorem}

\begin{proof}
  Use the sequential criterion (Theorem~\ref{thm:sequential-criterion-continuity}).\\
  Suppose $(x_n)$ in $U$ s.t. $x_n\to x_0$.
  \proofgoal{WTS:}{$(f+g)(x_n)\to(f+g)(x_0)$.}
  \[
    f(x_n)\to f(x_0),\qquad g(x_n)\to g(x_0).
  \]
  By the sequence limit laws (Theorem~\ref{thm:sequence-limit-laws}),
  \[
    f(x_n)+g(x_n)\to f(x_0)+g(x_0).
  \]
  The product law likewise gives $f(x_n)g(x_n)\to f(x_0)g(x_0)$.
  For the quotient, $g(x_0)\ne0$ allows the quotient law:
  $f(x_n)/g(x_n)\to f(x_0)/g(x_0)$ for sequences in the domain of $f/g$.
\end{proof}

\begin{theorem}{Continuity of Composition}{continuity-composition}
  Let $f:U\to\R$, $g:V\to\R$, for $U,V\subseteq\R$.
  $f(U)\subseteq V$, so $g\circ f$ is well-defined.
  \[
    f\text{ cont. at }x_0,\quad g\text{ cont. at }f(x_0)
    \implies g\circ f\text{ cont. at }x_0.
  \]
\end{theorem}

\begin{proof}
  By the sequential criterion (Theorem~\ref{thm:sequential-criterion-continuity}),
  fix $(x_n)$ in $U$ with $x_n\to x_0$.
  \proofgoal{WTS:}{$g(f(x_n))\to g(f(x_0))$.}
  Continuity of $f$ gives $f(x_n)\to f(x_0)$.
  Since $f(U)\subseteq V$, this is a sequence in $V$.
  Continuity of $g$ at $f(x_0)$ gives the desired limit.
\end{proof}
